\Needspace{9\baselineskip}
\noindent\SourceLeaf{16}\hspace*{\parindent}Drittes Beispiel.\qquad $p=13,\qquad q=11$
\[(3+2\sqrt{-1})^3=-9+46\sqrt{-1},\quad A\equiv-B\quad\text{Casus IV}\]
\[(-11)^3\equiv8\equiv\dfrac{3}{2.b}.\]
\Needspace{9\baselineskip}
\noindent\hspace*{\parindent}Viertes Beispiel.\qquad $p=17,\qquad q=11$
\[(1+4\sqrt{-1})^3=-47-52\sqrt{-1},\quad-47=A,\quad+52=B\]
\[A\equiv B\quad\text{Casus III}\]
\[(-11)^4\equiv4\equiv-\tfrac14.\]
\begin{center}[4.]\end{center}
\noindent\hspace*{\parindent}Es sei $p$ eine Primzahl, $n$ eine Primzahl $=3m+1$, $R$ Wurzel der Gleichung $x^3-1=0$, $r$ Wurzel der Gleichung $r^n-1=0$, $g$ Radix primitiva für mod. $n$.
\[r+Rr^g+R^2r^{g^2}+R^3r^{g^3}+\ldots+R^{n-2}r^{g^{n-2}}=P,\]
\[P^3=\tfrac12 n(M+N\sqrt{-27}),\]
\noindent wo
\[M\equiv1\ (\mathrm{mod.}\ 3),\qquad0\equiv M+3N(g^m-g^{2m})\ (\mathrm{mod.}\ n)\]
\[3k=M+2\]
\[a=\tfrac19(n+1+M)\]
\[b-c=N.\]
\[P^{pp}\equiv r^{pp}+Rr^{ppg}+\ldots\equiv P(R^{\mathrm{ind.}\,p})\]
\begingroup\footnotesize
\[\begin{array}{r@{}l|c}
&1&0\\[3pt]
\{\tfrac12 n(M+N\sqrt{-27})\}^{\frac13(pp-1)}\equiv{}&-\tfrac12+\tfrac12\sqrt{-3}\ (\mathrm{mod.}\ p)&\mathrm{ind.}\ p\equiv1\ (\mathrm{mod.}\ 3)\\[3pt]
&-\tfrac12-\tfrac12\sqrt{-3}&2
\end{array}\]
\endgroup
\[4n=(M+N\sqrt{-27})(M-N\sqrt{-27})\]
\[\dfrac{M}{3N}\equiv g^{2m}-g^m,\qquad-1\equiv g^{2m}+g^m\]
\begingroup\footnotesize
\[\begin{array}{r@{}l|c}
&1&1\\[3pt]
\left\{\dfrac{M-N\sqrt{-27}}{M+N\sqrt{-27}}\right\}^{\frac13(pp-1)}\equiv{}&-\tfrac12+\tfrac12\sqrt{-3}\ (\mathrm{mod.}\ p)&p^{\frac13(n-1)}\equiv-\tfrac12-\tfrac12\dfrac{M}{3N}\ (\mathrm{mod.}\ n)\\[5pt]
&-\tfrac12-\tfrac12\sqrt{-3}&-\tfrac12+\tfrac12\dfrac{M}{3N}
\end{array}\]
\endgroup
